\documentclass[11pt,a4paper,UTF8,fontset=none]{ctexart}
\setCJKmainfont{Noto Serif CJK SC}
\setCJKsansfont{Noto Sans CJK SC}
\setmainfont{DejaVu Serif}
\usepackage[a4paper,margin=22mm]{geometry}
\usepackage{amsmath,amssymb,booktabs,longtable,array,tabularx,xcolor,hyperref,fancyhdr,enumitem,microtype}
\hypersetup{colorlinks=true,linkcolor=blue,urlcolor=blue}
\definecolor{navy}{HTML}{173452}
\definecolor{sky}{HTML}{E9F4FA}
\pagestyle{fancy}\fancyhf{}\fancyhead[L]{\sffamily ZTQI · 研究试验版 v0.2}\fancyhead[R]{\sffamily 2026-10-21}\fancyfoot[C]{\thepage}
\setlength{\parindent}{2em}\setlength{\parskip}{4pt}
\begin{document}
\begin{center}
{\Huge\sffamily\bfseries ZTQI 做题区指数}\par\vspace{6mm}
{\LARGE\bfseries 构念建构、量表编制与效度检验方案}\par\vspace{4mm}
{\large ZTQI-72 研究试验版（60 核心条目 + 12 辅助条目）}\par\vspace{3mm}
{\normalsize ZTQI 研究项目 · 方法学方案 v0.2 · 2026 年 10 月 21 日}
\end{center}
\vspace{8mm}
\noindent\colorbox{sky}{\parbox{0.95\linewidth}{\textbf{重要声明}：这是一份研究设计与量表原型，不是已完成实证研究的论文。\textbf{没有招募数据、没有已验证的常模或信效度，也没有任何临床诊断效力。} ZTQI 不能用来公开排名、羞辱、自动判定某个人的道德或心理状况。}}
\begin{abstract}
本文提出 ZTQI（ZuoTiQu Index，做题区指数）的操作性构念：个体把学业成绩与解题能力泛化为自我价值、对他人的等级评价、人际回报资格、现实问题的标准答案偏好以及对非应试价值的贬抑的倾向。为避免把刻苦学习和优秀成绩误认为该构念，设计五维、每维 12 题共 60 条核心条目，并增设 12 条辅助的注意力、社会赞许与外部效标候选条目，形成共 72 题的 ZTQI-72 试验问卷。本文预先给出评分流程、反向计分、缺失值规则、结果披露边界、纸笔 OMR 流程与分阶段效度研究方案。全部评分权重为先验可修改参数，任何阈值均不得在尚无验证的条件下解释为标准化诊断界值。
\end{abstract}
\tableofcontents\newpage
\section{研究背景与构念边界}
“做题区”在网络语境中带有贬义且使用边界不稳定。ZTQI 不声称为网络语义提供唯一权威定义；相反，本研究先提出可证伪的操作性定义，再检验其可测性。研究对象是\textbf{态度倾向而非考试技能}。高分、勤奋、竞赛参与、社交内向、被拒绝或单次情绪波动均不是充分条件。
有关自我价值条件性的理论表明，个体可能依赖学业等特定领域维持自我价值\cite{crocker2003,grades2003}；学业应得感已有其他操作性定义\cite{chowning2009}。ZTQI 的 C 维强调\emph{对他人的人际回报期待}，与将学业成功归因于教师责任的学业应得感不同，不能简单视为既有量表的翻译或授权改编。
\section{理论模型与维度}
\begin{longtable}{p{1.2cm}p{4.9cm}p{2cm}p{5.1cm}}
\toprule 代码&构念维度&先验权重&排除边界\\\midrule
A&学业条件性自我价值&20\%&不等于高目标或正常考试焦虑\\
B&成绩等级化与人格泛化&30\%&不等于承认学业能力差异\\
C&优绩回报应得感&25\%&不等于希望自己的努力被看见\\
D&标准答案思维泛化&15\%&不等于逻辑严谨或解决问题能力\\
E&非应试价值贬低&10\%&不等于个人偏好学习活动\\\bottomrule
\end{longtable}
B、C 是用于识别网络标签常见权力与等级意涵的核心；A 是内在价值依赖；D、E 是外扩风格。五维模型是假设而非结论，若数据不支持高阶单因子，应只报告各维度，不强制求总指数。
\section{题库、作答格式与版本控制}
ZTQI-72-PILOT-0.2 共有 72 个条目，按固定随机顺序呈现：A--E 各 12 条正式计分题（20 条反向题），辅助部分 Q01--Q12 共 12 题。采用 5 级 Likert 反应：1 非常不同意、2 比较不同意、3 不确定/中立、4 比较同意、5 非常同意。全部题目源于原创编写而非复制既有受版权保护量表。问题编号和题目 ID 的映射由 questions.json 唯一规定，不能在纸质问卷、在线问卷和 OMR 识别中分别改写。
辅助题中 Q1--Q4 为社会赞许反应风格候选指标，Q5--Q6 为明确指定选择的注意力题，Q7--Q12 是拟用于外部效标的自陈经历条目；\textbf{辅助题不进入 ZTQI 总分}，且它们的测量意义本身待验证。注意力检查失败不能直接推断答题者不诚实。
\section{计分与缺失值处理}
记维度 $j\in\{A,B,C,D,E\}$ 的第 $i$ 项原始回答为 $x_{ji}\in\{1,2,3,4,5\}$。反向题转换为 $x^*_{ji}=6-x_{ji}$，正向题为 $x^*_{ji}=x_{ji}$。维度标准分
\begin{equation}
 D_j=25\left(\frac{1}{12}\sum_{i=1}^{12}x^*_{ji}-1\right),\qquad 0\le D_j\le100.
\end{equation}
若测量学验证支持可解释的综合指数，则暂定
\begin{equation}
 \mathrm{ZTQI}=0.20D_A+0.30D_B+0.25D_C+0.15D_D+0.10D_E.
\end{equation}
此权重不是统计估计，仅是试验性专家设定。\textbf{在试验版中，全部 60 个核心题完成才计算正式结果}；如果任何核心题缺失，显示“不完整，无法计分”，不做平均插补。辅助题缺失可提示数据质量但不改变总分。高低切点和百分位常模\textbf{暂不提供}，避免伪标准化。\textbf{结果必须同时给出五维原始分，禁止只展示单一标签。}
\section{示范性构成特征判别（不用于正式判定）}
历史草案中的 $A,B,C\ge60$ 和 $(0.20A+0.30B+0.25C)/0.75\ge70$ 只能作为\textbf{预注册的待验证假设}，不得用于网络公示、筛选人员或判断特定同学是否属于某一污名化群体。研究应比较连续特征模型与潜在类别模型，验证是否真的存在离散分类；没有明确证据时默认不存在可用的类别界限。
\section{样本、分析与验证计划}
\subsection{认知访谈与内容效度}
首先由不同背景参与者进行认知访谈（建议约 20--30 人），检查同义反复、社会赞许、文化语境歧义和暗含“成绩好的人有问题”的诱导性表述。邀请方法学研究者以专家评定进行条目筛选。
\subsection{预测试与因子结构}
计划获得多样化自愿样本，研究设计阶段预估 300--500 份用于初步项目分析和探索性因子分析（EFA），以及独立的 300--500 份用于验证性因子分析（CFA）；所需样本应通过模型复杂度、共同度与模拟功效分析重估。对有序分类题优先采用多分序相关矩阵和适合有序题的估计方法。评估单因子、五因子相关模型、二阶因子与双因子模型，检验区分效度。
\subsection{信度和效标效度}
分维报告 McDonald's $\omega$、必要时的 $\alpha$ 及区间估计；于 2--4 周后在同意重测的子样本上评价重测稳定性。关联学业条件性自尊、学业应得感、学校投入与同伴尊重态度等相关/不相关效标。检验题项对不同性别、年级及考试成绩群体的差异项功能（DIF），防止将分数优势错判为偏见。
\subsection{评分模型再估计与稳健性}
在验证集比较等权重、先验权重、潜变量分数模型。报告因子载荷、模型拟合、缺失比例、题项信息及稳健性检验。若综合分数的实质意义不足，撤销总分只保留维度画像。
\section{数据治理与伦理}
被试原则上自愿参与，可中途退出且无需解释；未成年参与研究需按适用法规和机构伦理要求安排知情同意。在线站点的默认设计是\textbf{浏览器本地计算、无上传、无第三方分析 SDK}；纸质答题卡扫描也只在浏览器 Canvas 中处理。研究数据的集中采集需另行经过伦理审查、明确目的、保留期限、撤回机制、访问控制与数据最小化设计。
本量表不得用于同学排名、同伴匿名指控、教师纪律处分、招生就业、心理诊断或其他高风险个人决策。结果是\textbf{自陈题项的研究性描述}，不是一个人的人格价值分数。
\section{答题卡自动识别协议}
纸张 A4，黑白印刷，建议 100\% 原尺寸打印。表单含四个黑色定位方块，中心坐标依次为 $ (12,12),(198,12),(198,285),(12,285)$ mm，坐标以页面左上为原点。72 行分左右各 36 行，填涂中心为
\begin{align}
 y_r &= 50+6.3(r\bmod36)\ \text{mm},\\
 x_{r,k} &= \begin{cases}41+8k,&0\le r<36,\\134+8k,&36\le r<72,\end{cases}\ \text{mm},\quad k=0,\ldots,4.
\end{align}
其中 $r$ 从零计数，第 $r+1$ 题的圆心由上式定位。定位器通过用户依次点击四角定位块中心，计算单应性投影并对每行五个泡泡取中心区域填涂密度。低置信度、空选和多选必须标记以人工复核，\textbf{不得无声自动补选}。识别正确性应由带真实不同扫描设备的标注样本评估，尚未完成该验证。
\section{局限与阶段性结论}
当前模型存在概念依赖于特定中文网络语境、构念间重叠、反向题方法因子、答题呈现方式效应与真实同伴行为差异等局限。\textbf{本稿仅交付可复现的研究方案和工具，绝不宣称发现了某种人群的实际比例或相关系数。}
\begin{thebibliography}{9}
\bibitem{crocker2003} Crocker, J., Luhtanen, R. K., Cooper, M. L., \& Bouvrette, A. (2003). Contingencies of self-worth in college students: Theory and measurement. \emph{Journal of Personality and Social Psychology}, 85(5), 894--908. DOI: 10.1037/0022-3514.85.5.894.
\bibitem{grades2003} Crocker, J., Karpinski, A., Quinn, D. M., \& Chase, S. K. (2003). When grades determine self-worth. \emph{Journal of Personality and Social Psychology}, 85(3), 507--516. DOI: 10.1037/0022-3514.85.3.507.
\bibitem{chowning2009} Chowning, K., \& Campbell, N. J. (2009). Development and validation of a measure of academic entitlement. \emph{Journal of Educational Psychology}, 101(4), 982--997. DOI: 10.1037/a0016351.
\bibitem{standards2014} AERA, APA, \& NCME. (2014). \emph{Standards for Educational and Psychological Testing}. AERA.
\end{thebibliography}
\appendix
\section{完整题库清单（固定纸笔顺序）}
\begin{longtable}{p{0.9cm}p{1.2cm}p{11.1cm}}
\toprule 编号&条目 ID&题干\\\midrule\endhead
01&C01&我投入很多时间学习后，会觉得别人理应给予我额外的尊重\\
02&Q1&我从来没有因为嫉妒而说过任何违心的话\\
03&B04&成绩较差的人也可能在重要问题上比我判断得更准确\\
04&D10&我愿意接受一些重要的人生问题没有最优解\\
05&D03&别人意见与我不同时，我常把分歧看作有人没有理解标准答案\\
06&B09&我觉得学业成绩可以作为评价一个人整体能力的主要标准\\
07&Q11&过去一个月，我参与过与考试成绩没有直接关系的兴趣活动\\
08&C05&我付出的学习努力应当换来同学更多的关注\\
09&A09&我担心成绩平庸会使我变成一个没有价值的人\\
10&Q9&过去一个月，我曾主动了解成绩不如我的同学擅长的领域\\
11&C09&我有时觉得学业优秀应当抵消我在人际交往中的不足\\
12&C10&被拒绝时，我能够尊重对方的自主选择而不以成绩辩解\\
13&C12&我不认为成绩优势能赋予我更多的交友或恋爱资格\\
14&C08&当自己的优秀没有得到认可时，我容易认为是别人不识货\\
15&C07&努力学习并不能使别人负有欣赏我的义务\\
16&B02&成绩不好的同学提出观点时，我会先怀疑其判断能力\\
17&D11&对于价值观冲突，我常试图证明一方必然比另一方正确\\
18&C11&当成绩不如我的人更受欢迎时，我会觉得秩序颠倒了\\
19&B01&只知道两个人成绩时，我会倾向于认为成绩更好的人整体上更优秀\\
20&B10&一个人的成绩与他应得的基本尊重没有关系\\
21&D12&理解不同人的经历，有时比给出答案更重要\\
22&B07&我不会根据一次考试成绩推断一个人的品格\\
23&C04&别人有权不喜欢我，即使我的成绩比大多数人更好\\
24&A03&我的自信很大程度上来自考试排名\\
25&E12&不同领域的努力都可能产生值得欣赏的结果\\
26&Q3&我从未有过任何不公平的想法\\
27&D06&只要掌握足够多的规则，我觉得人际交往就能像解题一样稳定\\
28&E03&如果一个人的成绩差，我会觉得他在其他领域的成功价值有限\\
29&A05&如果未来不再擅长考试，我会觉得自己的身份失去了支撑\\
30&E07&一个不擅长考试的人，也完全可能获得有意义的成就\\
31&D01&遇到人际冲突时，我首先想找出唯一正确的一方\\
32&B06&成绩比我差很多的人对我的建议，我通常不会太在意\\
33&B12&我能真诚欣赏成绩远不如我的人的其他长处\\
34&Q5&为核对问卷填写质量，请在本题选择“比较不同意”（2）\\
35&B11&当一个人成绩好时，我会更容易原谅他在人际相处中的缺点\\
36&A11&即使其他方面进步很大，考试失利仍会让我觉得自己一无是处\\
37&D09&我常把相处中的挫折归因于自己没找到正确解法\\
38&E08&若无法提升成绩，我认为多数课外活动不值得投入\\
39&E06&相比于考试高分，我通常不太看重一个人的创造力\\
40&E05&我认为将大量精力放在非应试活动上，大多是在浪费潜力\\
41&Q4&我永远不会为了表现好而隐瞒自己的真实想法\\
42&D07&即使无法明确判定对错，讨论仍然可能富有价值\\
43&A10&我是否值得被尊重，并不取决于最近一次考试\\
44&A12&我可以接受自己在某些学科上并不出色\\
45&B05&在评价一个陌生人时，我通常把学习成绩放在最重要的位置\\
46&E09&我很难理解有人自愿选择不以学业竞争为中心的生活\\
47&B08&在学校里，成绩越优秀的人一般越有资格获得话语权\\
48&A07&没有优异成绩，我也能坦然地向他人介绍自己\\
49&A06&我经常用分数来判断自己今天过得是否有意义\\
50&C03&我认为学业表现好的人在人际关系中理应有某些优先权\\
51&Q6&为核对问卷填写质量，请在本题选择“比较同意”（4）\\
52&A08&别人称赞我成绩好时，我才比较确信自己足够优秀\\
53&E10&尊重他人的生活选择并不要求其符合我的升学目标\\
54&Q7&过去一个月，我曾因考试成绩影响对某位同学的态度\\
55&D02&我认为大多数复杂的人际问题都可以归结为一套正确规则\\
56&D04&人际分歧有时确实可以同时存在多个合理立场\\
57&D08&复杂情境中，我通常不喜欢接受含糊不清的结论\\
58&A04&即使考试失利，我仍能清楚说出自己值得肯定的特点\\
59&Q12&过去一个月，我曾与他人讨论没有唯一答案的价值问题\\
60&B03&考试排名往往可以反映一个人是否值得敬佩\\
61&E11&在我看来，学习成绩一般比诚实、善良等品质更能说明人的价值\\
62&A01&当成绩下降时，我很难保持对自己的基本认可\\
63&C02&如果我成绩突出却没有得到某个人的好感，我会觉得不太公平\\
64&Q8&过去一个月，我曾因考试失利回避与人交往\\
65&Q2&我在所有情境下都会完全客观地评价别人\\
66&C06&同等条件下，我认为成绩更好的人更应得到朋友的偏爱\\
67&D05&面对社会问题，我经常只关注是否能找到明确的对错判断\\
68&E04&没有升学优势的兴趣爱好也可以是人生的重要组成部分\\
69&E01&我认为艺术或体育方面的成就通常不如高考试成绩值得尊重\\
70&A02&一次重要考试不理想会让我怀疑自己作为一个人的价值\\
71&E02&花很多时间发展与升学无关的兴趣，通常是不明智的\\
72&Q10&过去一个月，我曾因为他人的拒绝而联想到自己的成绩\\
\bottomrule\end{longtable}
\section{机器可读的题库与答题卡版本}
机器可读配置文件 \texttt{questions.json} 中固定存储 \texttt{version, number, id, dimension, reverse, scored, expected} 等字段。发生任何条目改动，应更新版本并同步更新纸质版、在线版、识别器与分析器，旧数据不允许无记录自动重计分。\end{document}